% Pregunta 204 (2021-julio-OEP2020, n.º 9): a la izquierda, el monopolio simple produce Q_m < Q_c (pérdida de eficiencia), salvo con discriminación de primer grado; a la derecha, su demanda de trabajo IPMg_L queda bajo el VPMg_L y contrata L_m < L_c con el salario competitivo. Autor: Víctor Gutiérrez Marcos
\documentclass[tikz,border=6pt]{standalone}
\input{tikz-tcee}
\begin{document}
\begin{tikzpicture}[tcee,x=1cm,y=1cm]
  % ---------- Panel 1: mercado del bien X ----------
  % D: P = 4,5 - 0,6Q ; IMg: P = 4,5 - 1,2Q ; CMg: P = 0,9 + 0,6Q
  % Monopolio: Qm = 2, Pm = 3,3, CMg(Qm) = 2,1 ; competencia: Qc = 3, Pc = 2,7
  \node[font=\small\bfseries] at (3,5.4) {Mercado del bien $X$};
  \fill[area] (2,3.3) -- (2,2.1) -- (3,2.7) -- cycle;
  \ejes{$Q$}{$P$}{6}{5}
  \draw[curva1] (0.3,4.32) -- (5.6,1.14) node[etiqueta,right] {$D$};
  \draw[curva2] (0.3,4.14) -- (3.6,0.18) node[etiqueta,right] {$IMg$};
  \draw[curva4] (0.3,1.08) -- (5.6,4.26) node[etiqueta,right] {$CMg$};
  \draw[guia] (2,0) node[below] {$Q_m$} -- (2,3.3) -- (0,3.3) node[left] {$P_m$};
  \draw[guia] (3,0) node[below] {$Q_c$} -- (3,2.7) -- (0,2.7) node[left] {$P_c$};
  \node[punto] at (2,3.3) {};
  \node[punto] at (2,2.1) {};
  \node[punto] at (3,2.7) {};
  \draw[->,thin,tceegris] (3.4,3.7) node[right,font=\footnotesize,text=black,align=left] {pérdida de\\eficiencia} -- (2.45,2.75);
  \node[anchor=base west,font=\footnotesize,align=left] at (-0.2,-1.2) {Discriminación de 1.\textsuperscript{er} grado: $IMg = D$,\\produce $Q_c$ y no hay pérdida (falla el «siempre» de la b).};

  % ---------- Panel 2: mercado de trabajo ----------
  % VPMg_L: w = 4,4 - 0,6L ; IPMg_L: w = 3,8 - 0,9L ; w = 2
  % Lm = 2 ; Lc = 4
  \begin{scope}[xshift=8.2cm]
  \node[font=\small\bfseries] at (3,5.4) {Mercado de trabajo};
  \ejes{$L$}{$w$}{6}{5}
  \draw[curva1] (0.3,4.22) -- (5.6,1.04) node[etiqueta,right] {$VPMg_L$};
  \draw[curva2] (0.3,3.53) -- (3.9,0.29) node[etiqueta,right] {$IPMg_L$};
  \draw[curva4] (0,2) -- (5.6,2);
  \node[left] at (0,2) {$w$};
  \draw[guia] (2,0) node[below] {$L_m$} -- (2,2);
  \draw[guia] (4,0) node[below] {$L_c$} -- (4,2);
  \node[punto] at (2,2) {};
  \node[punto] at (4,2) {};
  \node[anchor=base west,font=\footnotesize,align=left] at (-0.2,-1.2) {Discriminación de 1.\textsuperscript{er} grado: $IPMg_L = VPMg_L$,\\y contrata $L_c$, como la industria competitiva.};
  \end{scope}
\end{tikzpicture}
\end{document}
